\section{第三方库模板与公开后端契约}
\key{docs/libraries/starter/} 提供独立的第三方库模板，包含仅使用公开接口的源码、示例、三引擎检查器和兼容说明。
\begin{codeonly}
\usepackage[libraries=none]{gradbars}
\gradbarsdeclarelibrary{starter}{table}{gradbars-contrib-starter.sty}
\usegradbarslibrary{starter}
\end{codeonly}
请把第三方文件放在 TeX 搜索路径内。复制模板后统一修改 starter 的宏包、图形、格式化器、样式与后端名称，
在注册时声明依赖，并在库内检查依赖。
\key{\gradbarsversion} 返回宏包版本，\key{\gradbarsrequireversion{1.1.0}} 检查最低语义版本。
模板的独立检查入口为：
\begin{codeonly}
texlua docs/libraries/starter/check-starter.lua --gradbars-root .
\end{codeonly}
\subsection{标量后端的公开契约}
\key{\gradbarsbackendversion} 当前为 schema 1，为条形与 Heat 回调提供局部只读记录，私有模型存储仍不作为扩展接口。
\begin{codeonly}
\gradbarsdeclarebackend{my-backend}
 {bar,heat,missing,signed,gradient,print,labels,units,clip,
  decorations,embedded,tagging,artifact}
 {\gradbarsbackenddefault}
\gradbar[backend=my-backend]{72.125}
\end{codeonly}
回调的第一个参数为 bar 或 heat；只能声明实际实现或委托实现的能力。
\key{\gradbarsbackendvalue{字段}} 在回调期间读取未舍入数值、归一化位置、物理尺寸、选项与描述；
\key{\gradbarsbackendlabel} 排版原有格式化标签，\key{\gradbarsbackenddefault} 至多绘制一次捕获的原生续体。
\key{\gradbarsifbackendcapable{名称}{能力}{是}{否}} 查询声明。
未知后端、能力不足、递归回调及回调之外读取记录均有诊断。
复合图形与图例需显式选择 \key{backend fallback=tikz}，默认报错以免静默改变输出。
默认绘制使用 TikZ；公开后端契约不提供另一套内置绘图实现。
表格试排可能多次调用回调，因此不要在回调中写外部文件或产生其他可见副作用；
应尊重暂停标记状态，保留原始数值、单位、裁剪与基线。完整字段和能力说明见 \key{docs/backends.md}。

\clearpage
\subsection{独立的数值对注册与只读记录}
\label{sec:pair-extensions-cn}
比较两个观测时，可以为参考值与当前值建立独立的图形类型。
\key{\gradbarsdeclarepairtype} 的三个参数依次为类型名称、原生前端和回调；
前端选择 compare、dumbbell 或 range。
\key{\gradpairgraphic} 按“选项、类型、左值、右值”调用。
它不会把复合数据传给旧标量回调，原有 backend schema 1 保持不变。

\begin{example}[text and listing,sidebyside=false]{注册一个保留原生绘制的数值对类型}
\gradbarsdeclarepairtype{manual-pair}{dumbbell}{\gradbarspairdefault}
\gradpairgraphic[min=0,max=100,width=35mm,precision=1]
  {manual-pair}{66.125}{82.875}
\end{example}
compare 与 dumbbell 的输入角色是参考值、当前值；原生标签的顺序则为当前值/参考值。
range 的两个输入是下界、上界，不自动解释为误差或置信区间。
比较类允许任意一侧或两侧缺失；区间仍要求两端同时有效或同时缺失。

\subsubsection{在回调内读取原值与缺失状态}
回调接收的两个参数是规范原值或 NA，未经过显示精度舍入。
\key{\gradbarspairvalue} 可以读取角色、缺失、尺度、位置、尺寸与描述；
回调外读取或请求未知字段都会报错。
例如下面在原图之后附加原值之差；不完整的一对显示 n/a。
\begin{example}[text and listing,sidebyside=false]{原值之差与显示舍入分开处理}
\ExplSyntaxOn
\gradbarsdeclarepairtype{manual-pair-delta}{dumbbell}
 {\gradbarspairdefault\quad
  \str_if_eq:eeTF{\gradbarspairvalue{complete}}{true}
   {\gradbarsformatnumber[precision=1]{\gradbarspairvalue{delta}}}
   {n/a}}
\ExplSyntaxOff
\gradpairgraphic[max=100,width=30mm,precision=1]
  {manual-pair-delta}{66.125}{82.875}\par\medskip
\gradpairgraphic[max=100,width=30mm,precision=1]
  {manual-pair-delta}{NA}{0}
\end{example}
第一行先用原值求出差 16.75，再显示为一位小数；
第二行的零是有效观测，但参考值缺失，不能编造一个差值。

\begin{center}\small
\begin{tabularx}{\linewidth}{@{}l X@{}}
\toprule
\textbf{主要字段} & \textbf{含义}\\\midrule
\key{left}, \key{right} & 未舍入的规范数值，缺失侧为 NA。\\
\key{left-role}, \key{right-role} & reference/current 或 lower/upper。\\
\key{missing}, \key{complete} & none、left、right、both；两侧都有效时 complete 为 true。\\
\key{min}, \key{max}, \key{zero} & 已解析的公共范围，以及归一化后的零位置。\\
\key{left-position}, \key{right-position} & 裁剪到 0--1 的位置；缺失侧为 NA。\\
\key{left-clipped}, \key{right-clipped} & 原值是否超出显示范围。\\
\key{delta} & 右值减左值；不完整时为 NA，区间中表示端点跨度。\\
\key{width}, \key{height}, \key{alt} & 轨道物理尺寸与完整原生描述。\\
\bottomrule
\end{tabularx}\end{center}
裁剪只限制几何位置，不改原值或 delta；显示精度同样不参与几何和原值差计算。
完整字段、单位选项与接口签名见 \key{docs/pair-extensions.md}。

\subsubsection{能力查询与回调生命周期}
\key{\gradbarsifpairtypecapable} 接受类型、能力、满足时分支与不满足时分支。
例如 range 的 partial-missing 查询为否，dumbbell 为是。
未知类型或未知能力也走否分支。
\begin{codeonly}
\gradbarsifpairtypecapable{range}{partial-missing}
  {允许单侧缺失}{两端必须同时有效或同时缺失}
\end{codeonly}
\key{\gradbarspairdefault} 只绘制一次已捕获的原生续体，不再次调用公开前端；
即使放在额外盒子中，也不能重复委托。
\key{\gradbarspairlabel} 只排版已有双值标签，不负责位置或 Figure 标记。
递归调用数值对入口会报错；回调返回后记录访问立即关闭。
复用原生绘制可保留原始尺寸、基线、裁剪及替代文本；自行增加的文字和布局由回调负责。
表格试排可能重复执行回调，因此不得依赖外部写入、计数器或全局副作用。

\subsection{在原生表格中显式声明标准双值列}
\label{sec:pair-columns-cn}
独立数值对类型不会自动成为表格列。
\key{\gradbarsdeclarepaircolumn} 提供标准代理，显式选择 compare、dumbbell 或 range，
不接收任意独立消费者的排版回调。
声明采用普通 TeX 局部作用域，与原有列共享名称空间，重名会报错。
\begin{example}[text and listing,sidebyside=false]{配对报告：当前值、参考值与给定区间}
\begingroup
\gradbarsdeclarepaircolumn{manual-service}{frontend=dumbbell,
  options={min=0,max=100,precision=1,unit={\%}}}
\gradbarsdeclarepaircolumn{manual-span}{frontend=range,
  options={min=0,max=100,precision=1,unit={\%}}}
\begin{center}
\begin{gradtable}[style=paired-report,width=\linewidth,label width=25mm]
  \gradcolumn[width=25mm]{name}{服务项目}
  \gradcolumn[type=manual-service]{pair}{当前 / 参考}
  \gradcolumn[type=manual-span]{span}{给定区间}
  \gradrow{预约办理,66/82,60/88}
  \gradrow{上门服务,0/12,0/20}
  \gradrow{咨询服务,84/NA,NA}
\end{gradtable}
\end{center}
\endgroup
\end{example}
双值单元格按左值/右值填写，整格 NA 表示两端缺失。
范围与列宽沿用原生表格规则：自动范围收集两侧有效观测，
共享尺度显式绑定，最终分配的宽度仍优先。
配对列不参与标量筛选、排序、排名或聚合，也不隐式规范地区数字。
普通 tabular 与 tabularray 可在单元格中直接调用 \key{\gradpairgraphic}；
tabularray 应给足显式列宽，结构标记另按对应章节启用。

独立完整案例为 \key{docs/examples/paired-report-cn.tex}，
使用 \key{paired-report} 样式演示 accessible 和 mono 两种主题；
配套公开接口库为 \key{docs/libraries/gradbars-contrib-pairs.sty}。
从仓库根目录编译时，把 \key{sty/} 和 \key{docs/libraries/} 都加入 TeX 搜索路径。
接口准入及其他宿主限制见 \key{docs/pair-extensions.md}。
